\documentclass{article}
\usepackage[utf8]{inputenc}
\usepackage{booktabs}
\usepackage{array}
\usepackage{longtable}
\usepackage[landscape]{geometry}

\begin{document}

\title{Supplementary Material B: Per-study Quality Appraisal Scores}

\section*{Quality Appraisal Criteria}

The quality assessment framework employed five criteria, each scoring from 0 (No) to 1 (Yes), with 0.5 for ``To some extent'':

\begin{enumerate}
\item \textbf{QA1}: Are the objectives and the context of research clear, and well connected with DKT?
\item \textbf{QA2}: Does the research adequately delineate the methodology (e.g., training dataset, the model architecture, and its hyper-parameters)?
\item \textbf{QA3}: Does the research adequately delineate the experiments (e.g., examining sequential or temporal consistency, assessing interpretability, addressing data inconsistencies, or evaluating the model's real-world applicability)?
\item \textbf{QA4}: Is there congruence between methodology and methods used for data collection, analysis, and interpretation?
\item \textbf{QA5}: Are the study contributions and limitations clearly stated?
\end{enumerate}

Studies scoring $\geq 3.5$ points were classified as high quality, those between 1.5--3.5 points as medium quality, and those $\leq 1.5$ points as low quality.

\section*{Distribution Summary}

Total studies analyzed: 84

Quality distribution:
\begin{itemize}
\item High quality ($\geq 3.5$): 83 studies (98.8\%)
\item Medium quality (1.5--3.5): 1 study (1.2\%)  
\item Low quality ($\leq 1.5$): 0 studies (0.0\%)
\end{itemize}

Mean quality score: 4.32 (SD = 0.52)

\section*{Per-study Quality Appraisal Scores}

\begin{longtable}{p{1cm}p{8cm}p{1cm}p{1cm}p{1cm}p{1cm}p{1cm}p{1.5cm}p{1.5cm}}
\toprule
\textbf{ID} & \textbf{Study} & \textbf{QA1} & \textbf{QA2} & \textbf{QA3} & \textbf{QA4} & \textbf{QA5} & \textbf{Total} & \textbf{Quality} \\
\midrule
\endfirsthead

\multicolumn{9}{c}%
{\tablename\ \thetable\ -- \textit{Continued from previous page}} \\
\toprule
\textbf{ID} & \textbf{Study} & \textbf{QA1} & \textbf{QA2} & \textbf{QA3} & \textbf{QA4} & \textbf{QA5} & \textbf{Total} & \textbf{Quality} \\
\midrule
\endhead

\midrule \multicolumn{9}{r}{\textit{Continued on next page}} \\
\endfoot

\endlastfoot

1 & Calibrated Q-Matrix-Enhanced Deep Knowledge Tracing with Relational Attention Mechanism & 1.0 & 1.0 & 1.0 & 1.0 & 1.0 & 5.0 & High \\
2 & Graph-based knowledge tracing: Modeling student proficiency using graph neural networks & 1.0 & 1.0 & 1.0 & 1.0 & 1.0 & 5.0 & High \\
3 & Learning consistent representations with temporal and causal enhancement for knowledge tracing & 1.0 & 1.0 & 1.0 & 1.0 & 1.0 & 5.0 & High \\
4 & Knowledge Graph and Personalized Answer Sequences for Programming Knowledge Tracing & 1.0 & 1.0 & 1.0 & 1.0 & 1.0 & 5.0 & High \\
5 & Bridging the Vocabulary Gap: Using Side Information for Deep Knowledge Tracing & 1.0 & 1.0 & 1.0 & 1.0 & 1.0 & 5.0 & High \\
6 & ELAKT: Enhancing Locality for Attentive Knowledge Tracing & 1.0 & 1.0 & 1.0 & 1.0 & 1.0 & 5.0 & High \\
7 & Deep Knowledge Tracing with Learning Curves & 1.0 & 1.0 & 1.0 & 1.0 & 1.0 & 5.0 & High \\
8 & Design and evaluation of Trustworthy Knowledge Tracing Model for Intelligent Tutoring System & 1.0 & 1.0 & 1.0 & 1.0 & 1.0 & 5.0 & High \\
9 & Pull together: Option-weighting-enhanced mixture-of-experts knowledge tracing & 1.0 & 1.0 & 1.0 & 1.0 & 1.0 & 5.0 & High \\
10 & Heterogeneous graph-based knowledge tracing with spatiotemporal evolution & 1.0 & 1.0 & 1.0 & 1.0 & 1.0 & 5.0 & High \\
11 & Deep Knowledge Tracking Integrating Programming Exercise Difficulty and Forgetting Factors & 1.0 & 1.0 & 1.0 & 1.0 & 1.0 & 5.0 & High \\
12 & Temporal learner modelling through integration of neural and symbolic architectures & 1.0 & 1.0 & 1.0 & 1.0 & 1.0 & 5.0 & High \\
13 & A Knowledge Query Network Model Based on Rasch Model Embedding for Personalized Online Learning & 1.0 & 1.0 & 1.0 & 1.0 & 1.0 & 5.0 & High \\
14 & Interpretable Knowledge Tracing with Multiscale State Representation & 1.0 & 1.0 & 1.0 & 1.0 & 1.0 & 5.0 & High \\
15 & EKT: Exercise-Aware Knowledge Tracing for Student Performance Prediction & 1.0 & 1.0 & 1.0 & 1.0 & 0.5 & 4.5 & High \\
16 & Knowledge tracing machines: Factorization machines for knowledge tracing & 1.0 & 1.0 & 1.0 & 1.0 & 0.5 & 4.5 & High \\
17 & DF-EGM: Personalised Knowledge Tracing with Dynamic Forgetting and Enhanced Gain Mechanisms & 1.0 & 1.0 & 1.0 & 1.0 & 0.5 & 4.5 & High \\
18 & Context-Aware Attentive Knowledge Tracing & 1.0 & 1.0 & 1.0 & 1.0 & 0.5 & 4.5 & High \\
19 & Tracking knowledge proficiency of students with calibrated Q-matrix & 1.0 & 1.0 & 0.5 & 1.0 & 1.0 & 4.5 & High \\
20 & GameDKT: Deep knowledge tracing in educational games & 1.0 & 1.0 & 0.5 & 1.0 & 1.0 & 4.5 & High \\
21 & Knowledge interaction enhanced sequential modeling for interpretable learner knowledge diagnosis in intelligent tutoring systems & 1.0 & 1.0 & 0.5 & 1.0 & 1.0 & 4.5 & High \\
22 & Deep Graph Memory Networks for Forgetting-Robust Knowledge Tracing & 1.0 & 1.0 & 0.5 & 1.0 & 1.0 & 4.5 & High \\
23 & Ensemble Knowledge Tracing: Modeling interactions in learning process & 1.0 & 1.0 & 1.0 & 1.0 & 0.5 & 4.5 & High \\
24 & Federated Deep Knowledge Tracing & 1.0 & 1.0 & 1.0 & 1.0 & 0.5 & 4.5 & High \\
25 & Deep Learning vs. Bayesian Knowledge Tracing: Student Models for Interventions & 1.0 & 1.0 & 0.5 & 1.0 & 1.0 & 4.5 & High \\
26 & A novel framework for deep knowledge tracing via gating-controlled forgetting and learning mechanisms & 1.0 & 1.0 & 1.0 & 1.0 & 0.5 & 4.5 & High \\
27 & Infusing Expert Knowledge Into a Deep Neural Network Using Attention Mechanism for Personalized Learning Environments & 1.0 & 1.0 & 1.0 & 1.0 & 0.5 & 4.5 & High \\
28 & Exploiting multiple question factors for knowledge tracing & 1.0 & 1.0 & 1.0 & 1.0 & 0.5 & 4.5 & High \\
29 & AdaptKT: A Domain Adaptable Method for Knowledge Tracing & 1.0 & 1.0 & 1.0 & 1.0 & 0.5 & 4.5 & High \\
30 & Dynamic Cognitive Diagnosis: An Educational Priors-Enhanced Deep Knowledge Tracing Perspective & 1.0 & 1.0 & 0.5 & 1.0 & 1.0 & 4.5 & High \\
31 & Modeling knowledge proficiency using multi-hierarchical capsule graph neural network & 1.0 & 1.0 & 1.0 & 1.0 & 0.5 & 4.5 & High \\
32 & Monitoring Student Progress for Learning Process-Consistent Knowledge Tracing & 1.0 & 1.0 & 1.0 & 1.0 & 0.5 & 4.5 & High \\
33 & SPAKT: A Self-Supervised Pre-TrAining Method for Knowledge Tracing & 1.0 & 1.0 & 1.0 & 1.0 & 0.5 & 4.5 & High \\
34 & Adaptive meta-knowledge dictionary learning for incremental knowledge tracing & 1.0 & 1.0 & 0.5 & 1.0 & 1.0 & 4.5 & High \\
35 & GraphCA: Learning from Graph Counterfactual Augmentation for Knowledge Tracing & 1.0 & 1.0 & 0.5 & 1.0 & 1.0 & 4.5 & High \\
36 & Enhanced Knowledge Tracing With Learnable Filter & 1.0 & 1.0 & 1.0 & 1.0 & 0.5 & 4.5 & High \\
37 & Self-paced contrastive learning for knowledge tracing & 1.0 & 1.0 & 1.0 & 1.0 & 0.5 & 4.5 & High \\
38 & Global Feature-guided Knowledge Tracing & 1.0 & 1.0 & 1.0 & 1.0 & 0.5 & 4.5 & High \\
39 & MLC-DKT: A multi-layer context-aware deep knowledge tracing model & 1.0 & 1.0 & 1.0 & 1.0 & 0.5 & 4.5 & High \\
40 & Interaction Sequence Temporal Convolutional Based Knowledge Tracing & 1.0 & 1.0 & 1.0 & 1.0 & 0.5 & 4.5 & High \\
41 & Integrating fine-grained attention into multi-task learning for knowledge tracing & 1.0 & 1.0 & 1.0 & 1.0 & 0.5 & 4.5 & High \\
42 & A probabilistic generative model for tracking multi-knowledge concept mastery probability & 1.0 & 1.0 & 1.0 & 1.0 & 0.5 & 4.5 & High \\
43 & SFBKT: A Synthetically Forgetting Behavior Method for Knowledge Tracing & 1.0 & 1.0 & 1.0 & 1.0 & 0.5 & 4.5 & High \\
44 & Towards more accurate and interpretable model: Fusing multiple knowledge relations into deep knowledge tracing & 1.0 & 1.0 & 1.0 & 1.0 & 0.5 & 4.5 & High \\
45 & A Genetic Causal Explainer for Deep Knowledge Tracing & 1.0 & 1.0 & 1.0 & 1.0 & 0.5 & 4.5 & High \\
46 & Explore Bayesian analysis in Cognitive-aware Key-Value Memory Networks for knowledge tracing in online learning & 1.0 & 1.0 & 1.0 & 1.0 & 0.5 & 4.5 & High \\
47 & Meta-path structured graph pre-training for improving knowledge tracing in intelligent tutoring & 1.0 & 1.0 & 0.5 & 1.0 & 1.0 & 4.5 & High \\
48 & GELT: A graph embeddings based lite-transformer for knowledge tracing & 1.0 & 1.0 & 1.0 & 1.0 & 0.5 & 4.5 & High \\
49 & An efficient state-aware Coarse-Fine-Grained model for Knowledge Tracing & 1.0 & 1.0 & 1.0 & 1.0 & 0.5 & 4.5 & High \\
50 & Heterogeneous Evolution Network Embedding with Temporal Extension for Intelligent Tutoring Systems & 1.0 & 1.0 & 0.5 & 1.0 & 1.0 & 4.5 & High \\
51 & Leveraging Pedagogical Theories to Understand Student Learning Process with Graph-based Reasonable Knowledge Tracing & 1.0 & 1.0 & 0.5 & 1.0 & 1.0 & 4.5 & High \\
52 & JKT: A joint graph convolutional network based Deep Knowledge Tracing & 1.0 & 1.0 & 0.5 & 1.0 & 0.5 & 4.0 & High \\
53 & Learning Process-consistent Knowledge Tracing & 1.0 & 1.0 & 0.5 & 1.0 & 0.5 & 4.0 & High \\
54 & Implicit Heterogeneous Features Embedding in Deep Knowledge Tracing & 1.0 & 1.0 & 0.5 & 1.0 & 0.5 & 4.0 & High \\
55 & Knowledge Query Network for Knowledge Tracing & 1.0 & 1.0 & 0.5 & 1.0 & 0.5 & 4.0 & High \\
56 & SGKT: Session graph-based knowledge tracing for student performance prediction & 1.0 & 1.0 & 0.5 & 1.0 & 0.5 & 4.0 & High \\
57 & Knowledge structure enhanced graph representation learning model for attentive knowledge tracing & 1.0 & 1.0 & 0.5 & 1.0 & 0.5 & 4.0 & High \\
58 & Improving Knowledge Tracing with Collaborative Information & 1.0 & 1.0 & 0.5 & 1.0 & 0.5 & 4.0 & High \\
59 & Enhancing Knowledge Tracing via Adversarial Training & 1.0 & 1.0 & 0.5 & 1.0 & 0.5 & 4.0 & High \\
60 & Ability boosted knowledge tracing & 1.0 & 1.0 & 0.5 & 1.0 & 0.5 & 4.0 & High \\
61 & Assessing Student's Dynamic Knowledge State by Exploring the Question Difficulty Effect & 1.0 & 1.0 & 0.5 & 1.0 & 0.5 & 4.0 & High \\
62 & Enhancing Deep Knowledge Tracing with Auxiliary Tasks & 1.0 & 1.0 & 0.5 & 1.0 & 0.5 & 4.0 & High \\
63 & Introducing Problem Schema with Hierarchical Exercise Graph for Knowledge Tracing & 1.0 & 1.0 & 0.5 & 1.0 & 0.5 & 4.0 & High \\
64 & What is wrong with deep knowledge tracing? Attention-based knowledge tracing & 1.0 & 1.0 & 0.5 & 1.0 & 0.5 & 4.0 & High \\
65 & Learning Behavior-oriented Knowledge Tracing & 1.0 & 1.0 & 0.5 & 1.0 & 0.5 & 4.0 & High \\
66 & Knowledge tracing based on multi-feature fusion & 1.0 & 1.0 & 0.5 & 1.0 & 0.5 & 4.0 & High \\
67 & Improving knowledge tracing via a heterogeneous information network enhanced by student interactions & 1.0 & 1.0 & 0.5 & 1.0 & 0.5 & 4.0 & High \\
68 & DCKT: A Novel Dual-Centric Learning Model for Knowledge Tracing & 1.0 & 1.0 & 0.5 & 1.0 & 0.5 & 4.0 & High \\
69 & Transition-Aware Multi-Activity Knowledge Tracing & 1.0 & 1.0 & 0.5 & 1.0 & 0.5 & 4.0 & High \\
70 & Learning from Non-Assessed Resources: Deep Multi-Type Knowledge Tracing & 1.0 & 1.0 & 0.5 & 1.0 & 0.5 & 4.0 & High \\
71 & DKVMN-KAPS: Dynamic Key-Value Memory Networks Knowledge Tracing With Students' Knowledge-Absorption Ability and Problem-Solving Ability & 1.0 & 1.0 & 0.5 & 1.0 & 0.5 & 4.0 & High \\
72 & Discovering Multi-Relational Integration for Knowledge Tracing with Retentive Networks & 1.0 & 1.0 & 0.5 & 1.0 & 0.5 & 4.0 & High \\
73 & Improving Knowledge Tracing via Considering Students' Interaction Patterns & 1.0 & 1.0 & 0.5 & 1.0 & 0.5 & 4.0 & High \\
74 & Stable Knowledge Tracing Using Causal Inference & 1.0 & 1.0 & 0.5 & 1.0 & 0.5 & 4.0 & High \\
75 & Modeling Student Performance Using Feature Crosses Information for Knowledge Tracing & 1.0 & 0.5 & 1.0 & 1.0 & 0.5 & 4.0 & High \\
76 & A deeper knowledge tracking model integrating cognitive theory and learning behavior & 1.0 & 1.0 & 0.5 & 1.0 & 0.5 & 4.0 & High \\
77 & Attention and Learning Features-Enhanced Knowledge Tracing & 1.0 & 1.0 & 0.5 & 1.0 & 0.5 & 4.0 & High \\
78 & MonaCoBERT: Monotonic Attention Based ConvBERT for Knowledge Tracing & 1.0 & 1.0 & 0.5 & 1.0 & 0.5 & 4.0 & High \\
79 & Deep Knowledge Tracing Incorporating a Hypernetwork With Independent Student and Item Networks & 1.0 & 1.0 & 0.5 & 1.0 & 0.5 & 4.0 & High \\
80 & Embedding cognitive framework with self-attention for interpretable knowledge tracing & 1.0 & 1.0 & 0.5 & 1.0 & 0.5 & 4.0 & High \\
81 & Knowledge Tracing Model and Student Profile Based on Clustering-Neural-Network & 1.0 & 1.0 & 0.5 & 1.0 & 0.5 & 4.0 & High \\
82 & Improving Interpretability of Deep Sequential Knowledge Tracing Models with Question-centric Cognitive Representations & 1.0 & 1.0 & 0.5 & 1.0 & 0.5 & 4.0 & High \\
83 & Model-agnostic counterfactual reasoning for identifying and mitigating answer bias in knowledge tracing & 1.0 & 1.0 & 0.5 & 1.0 & 0.5 & 4.0 & High \\
84 & Enhanced Knowledge Tracing via Frequency Integration and Order Sensitivity & 1.0 & 1.0 & 0.5 & 1.0 & 0.5 & 4.0 & High \\

\botrule
\end{longtable}

\section*{Sensitivity Analysis}

Given that the vast majority of studies (98.8\%) achieved high quality scores and only one study was classified as medium quality, removing medium-quality studies would not substantially alter the key findings of this systematic review. The high mean quality score (4.32/5.0) indicates that methodological rigor in DKT research is generally strong across basic quality criteria.

\end{document}